\section{多候选联合AP与波形消元块NOMP}\label{sec:methods}
本章分别改进候选选择和局部更新。MCJ-AP保留多个条件起点，以联合精修后的残差择优；VP-BNOMP在试探位置时重新估计波形，减少固定系数带来的重复小步。两种方法使用同一位置目标，搜索和更新策略的影响可以据此比较。
\subsection{条件投影搜索}
寻找第二源时，应先区分哪些观测已由第一源解释。给定已发现位置，以$\bm B_f\in\C^{L\times K_{\mathrm B}}$收集其响应，$K_{\mathrm B}$为已固定的源数；搜索第二源时$K_{\mathrm B}=1$。对观测与候选响应作相同正交投影，定义
\begin{equation}
 \bm r_f=\bm\Pi_{\bm B_f}^{\perp}\bm x_f,\qquad
 \bm u_f(\bm p)=\bm\Pi_{\bm B_f}^{\perp}\bm a_f(\bm p).
\end{equation}
其中$\bm\Pi_{\bm B_f}^{\perp}=\bm I_L-\bm B_f\bm B_f^\dagger$，$\bm r_f\in\C^L$为投影后的观测，$\bm u_f(\bm p)\in\C^L$为候选响应经过同一投影后的向量。对应的条件位置谱为
\begin{equation}
 Q_{\mathrm{AP}}(\bm p\mid\bm B)
 =E_x^{-1}\sum_f\frac{|\bm u_f(\bm p)^H\bm r_f|^2}
 {\|\bm u_f(\bm p)\|_2^2}.\label{eq:conditional}
\end{equation}
条件谱衡量新增位置带来的拟合改善。将$\bm a_f$分解为已有响应空间内的部分与正交部分$\bm u_f$。当$\bm u_f\ne0$时，新方向与已有响应正交，投影矩阵满足
\begin{equation}
 \bm\Pi_{[\bm B_f,\bm a_f]}^{\perp}
 =\bm\Pi_{\bm B_f}^{\perp}
 -\frac{\bm u_f\bm u_f^H}{\|\bm u_f\|^2} .
\end{equation}
对上式作观测的二次型并求和，便有
\begin{equation}
 J_{\mathrm B}-J([\bm B,\bm a(\bm p)])
 =\frac1{E_x}\sum_f\frac{|\bm u_f(\bm p)^H\bm r_f|^2}
 {\|\bm u_f(\bm p)\|^2}
 =Q_{\mathrm{AP}}(\bm p\mid\bm B),\label{eq:conditionalGain}
\end{equation}
其中$J_{\mathrm B}=E_x^{-1}\sum_f\|\bm r_f\|^2$为仅拟合已有源的残差，$J([\bm B,\bm a(\bm p)])$表示加入候选后、共同重估所有系数所得的残差。若$\bm u_f=0$，该候选没有提供新的响应方向，此频点贡献记为零。

式\eqref{eq:conditionalGain}将条件搜索与联合拟合连接起来：固定已有位置时，条件得分最高的候选使联合残差最小\cite{ChenHudsonYao2002}。这一等价关系要求观测和候选响应使用相同投影。

\subsection{多候选联合精修AP（MCJ-AP）}\label{sec:mcj}
粗网格排序回答的是哪个网格点当前得分更高，连续精修则允许两源一起移动；二者得到的最佳分支可能不同。MCJ-AP据此保留多个可能起点。先在公共网格$\mathcal G\subset\Omega$上最大化$Q_0$，作单源连续精修，得到$\bm p_1^{(0)}$。固定该位置后，计算第二源条件谱，提取$3\times3$邻域内的局部极大值，再按得分从高到低保留间隔不少于$\Delta_{\mathrm c}$的点，得到$M$个候选$\bm c_1,\ldots,\bm c_M\in\R^2$。本文取$\Delta_{\mathrm c}=30$\,m。每个候选分别与第一源组成初值，进入相同的联合精修：
\begin{equation}
 \widehat{\bm P}^{(m)}=\operatorname{LocalMin}_{\bm P\in\Omega^2}
 J(\bm P;[\bm p_1^{(0)},\bm c_m]),\qquad
 m^*=\arg\min_m J(\widehat{\bm P}^{(m)}).\label{eq:branch}
\end{equation}
其中$\operatorname{LocalMin}$表示在搜索边界内从给定初值局部最小化$J$，$\widehat{\bm P}^{(m)}$为第$m$个分支的结果。算法输出最终残差最小的分支$\widehat{\bm P}^{(m^*)}$。因此，候选质量由精修后的完整拟合决定，粗谱仅负责提供起点。

\begin{algorithm}[H]
\caption{多候选联合精修AP}\label{alg:ap}
\begin{algorithmic}[1]
\Require 观测$\{\bm x_f\}$，站点坐标，公共网格$\mathcal G$，候选数$M$
\State 最大化$Q_0$并单源精修，得到$\bm p_1^{(0)}$
\State 固定$\bm p_1^{(0)}$，以式\eqref{eq:conditional}扫描$\mathcal G$
\State 提取邻域极大值，按条件得分和间距$\Delta_{\mathrm c}$保留$M$个候选
\For{$m=1,\ldots,M$}
  \State 从$[\bm p_1^{(0)},\bm c_m]$联合优化$J$；每次试探按式\eqref{eq:waveLS}重估波形
  \State 保存该分支的目标、位置与停止状态
\EndFor
\State 按式\eqref{eq:branch}输出最小残差分支
\end{algorithmic}
\end{algorithm}

联合精修采用有界拟牛顿法L-BFGS-B同时优化两个位置\cite{Byrd1995}。将坐标按源排列为$\bm\theta=\operatorname{vec}(\bm P)\in\R^{2K}$，其第$i$个分量为$\theta_i$；联合拟合残差为$\bm e_f=\bm x_f-\bm A_f\widehat{\bm s}_f\in\C^L$。响应矩阵局部满列秩时，最小二乘解满足$\bm A_f^H\bm e_f=0$。这一正交条件使求导中含$\partial\widehat{\bm s}_f/\partial\theta_i$的项消失，因而位置梯度可直接写为\cite{GolubPereyra1973}
\begin{equation}
 \frac{\partial J}{\partial\theta_i}
 =-\frac{2}{E_x}\Rea\sum_f
 \bm e_f^H\frac{\partial\bm A_f}{\partial\theta_i}
 \widehat{\bm s}_f,\label{eq:gradient}
\end{equation}
式中$\Rea$表示取实部，$\partial\bm A_f/\partial\theta_i\in\C^{L\times K}$为响应矩阵对该坐标的导数。记$\bm v_l=(\bm p-\bm r_l)/d_l\in\R^2$为从第$l$站指向候选位置的单位方向，$d_l=\|\bm p-\bm r_l\|_2$，响应的空间梯度为
\begin{equation}
 \nabla_{\bm p}a_l(f,\bm p)
 =a_l(f,\bm p)\left(-\frac1{d_l}-\mathrm j\frac{2\pi f}{c_0}\right)\bm v_l.
 \label{eq:responseDerivative}
\end{equation}
这里$a_l(f,\bm p)=[\bm a_f(\bm p)]_l$。括号内两项分别来自距离幅度和传播时延的变化。联合优化每次重估波形，但计算式\eqref{eq:gradient}时不需要显式求波形估计的导数。

增加候选数时，保持已有分支的初值、求解器与停止条件不变。较大集合包含较小集合的全部精修结果，故选择后的最小残差不会增大；对应的位置误差由实验评价。$M=1$为单候选版本，$M>1$以额外局部计算换取更多有效起点。

\subsection{波形消元块NOMP（VP-BNOMP）}\label{sec:vpnomp}
NOMP在当前拟合残差中发现新成分，再循环修正已有参数\cite{Mamandipoor2016}。二维多频率块NOMP将同一位置的逐频系数视为一个参数块\cite{Yang2022}。本文保留这一发现和循环结构，改进局部位置更新：每试探一个位置，就先求该源的最优信号系数，再判断位置是否更好。

给定当前位置和共同估计的系数$\widehat s_{k,f}$，更新源$j$时先减去其他源的当前拟合贡献：
\begin{equation}
 \bm y_{j,f}=\bm x_f-\sum_{k\ne j}\bm a_f(\bm p_k)\widehat s_{k,f} .
 \label{eq:targetResidual}
\end{equation}
其中$\bm y_{j,f}\in\C^L$为保留源$j$贡献的残差观测，在一次局部位置优化内固定。对任意试探位置$\bm p$，其最优复数系数$\widehat q_f(\bm p)\in\C$与局部目标$\phi_j(\bm p)$分别为
\begin{align}
 \widehat q_f(\bm p)&=\frac{\bm a_f(\bm p)^H\bm y_{j,f}}
 {\|\bm a_f(\bm p)\|^2},\label{eq:localcoef}\\
 \phi_j(\bm p)&=\min_{\{q_f\}}E_x^{-1}\sum_f
 \|\bm y_{j,f}-\bm a_f(\bm p)q_f\|^2\nonumber\\
 &=E_x^{-1}\sum_f\left(\|\bm y_{j,f}\|^2-
 \frac{|\bm a_f(\bm p)^H\bm y_{j,f}|^2}{\|\bm a_f(\bm p)\|^2}\right).
 \label{eq:localprofile}
\end{align}
这里$q_f$为标量信号系数，与单位范数响应向量$\bm q_f$区分。式\eqref{eq:localprofile}表示每个试探位置都配上其最优波形。Newton更新使用该消元目标的梯度与Hessian（二阶导数矩阵）；回溯时也重新求系数，仅接受局部残差不增的试探。更新后同步替换该源的响应、系数和显式残差。逐源循环结束后，再共同重估全部已发现源的系数。这样保持块NOMP的发现与循环结构，局部更新的区别集中在是否消去波形。

发现新源时，记当前共同拟合残差为$\bm r_f=\bm x_f-\sum_k\bm a_f(\bm p_k)\widehat s_{k,f}$；尚未发现源时$\bm r_f=\bm x_f$。扫描的残差谱为
\begin{equation}
 Q_{\mathrm R}(\bm p)=E_x^{-1}\sum_f
 \frac{|\bm a_f(\bm p)^H\bm r_f|^2}{\|\bm a_f(\bm p)\|^2},
 \label{eq:residualScore}
\end{equation}
该谱评价候选对当前残差的解释能力。它用于发现初始位置，后续再以联合拟合修正所有已发现源。

\begin{algorithm}[H]
\caption{波形消元块NOMP}\label{alg:nomp}
\begin{algorithmic}[1]
\Require 观测，公共网格，源数$K$，循环与Newton停止参数
\State 初始化源集合为空，$\bm r_f\gets\bm x_f$
\For{$k=1,\ldots,K$}
  \State 最大化残差谱\eqref{eq:residualScore}，发现新位置
  \State 以式\eqref{eq:localprofile}作单源Newton精修，加入源集合并更新残差
  \Repeat
    \For{已发现的每个源$j$}
      \State 构造$\bm y_{j,f}$，用$\phi_j$作有界Newton更新
      \State 每个试探位置重估$q_f$，仅接受局部残差不增的更新，同步替换该源贡献
    \EndFor
  \Until{循环停止条件满足或达到安全上限}
  \State 按式\eqref{eq:waveLS}共同重估全部已发现源的系数，更新共同拟合残差
\EndFor
\State 输出已发现位置
\end{algorithmic}
\end{algorithm}

两种发现方式的区别在于如何处理已有源。共同LS拟合后，残差$\bm r_f$与已有响应的列空间正交，因此AP条件谱和NOMP残差谱的分子相同；分母则分别为$\|\bm\Pi_{\bm B_f}^{\perp}\bm a_f\|^2$和$\|\bm a_f\|^2$。AP还考虑候选与已有响应的相似程度，NOMP直接扫描当前残差，两者因而可能选中不同的初始区域。本文保留这两种发现机制。

\subsection{消元更新的曲率与下降性质}\label{sec:curvature}
下面说明为什么固定信号系数可能使位置更新缓慢。对于一次源位置更新，将二维坐标记为$\bm z\in\R^2$，将$F$个复信号系数的实部和虚部排列为$\bm u\in\R^{2F}$。对应的归一化局部损失记为$\mathcal L_{\mathrm{loc}}(\bm z,\bm u)$，给定位置的最小二乘系数记为$\widehat{\bm u}(\bm z)$，消元后损失为$\varphi(\bm z)=\mathcal L_{\mathrm{loc}}(\bm z,\widehat{\bm u}(\bm z))$，即式\eqref{eq:localprofile}。

假设响应二阶连续可微，且系数最小二乘解局部唯一。定义$H_{zz}\in\R^{2\times2}$、$H_{uu}\in\R^{2F\times2F}$和$H_{zu}\in\R^{2\times2F}$分别为$\mathcal L_{\mathrm{loc}}$的位置、系数及混合二阶导数，$H_{uz}=H_{zu}^T$，所有导数在$(\bm z,\widehat{\bm u}(\bm z))$处计算。最优系数满足$\nabla_u\mathcal L_{\mathrm{loc}}=0$，对位置求导可得
\begin{equation}
 H_{uz}+H_{uu}\frac{\partial\widehat{\bm u}}{\partial\bm z}=0,
 \qquad\frac{\partial\widehat{\bm u}}{\partial\bm z}
 =-H_{uu}^{-1}H_{uz}.\label{eq:coefficientDerivative}
\end{equation}
将这一关系代入$\varphi$的一、二阶导数，得到变量投影的梯度和曲率关系\cite{GolubPereyra1973}：
\begin{align}
 \nabla\varphi&=\nabla_z\mathcal L_{\mathrm{loc}},\label{eq:envelope}\\
 H_{\mathrm{VP}}&=H_{zz}-H_{zu}H_{uu}^{-1}H_{uz},\label{eq:schur}\\
 H_{zz}-H_{\mathrm{VP}}&\succeq0 .\label{eq:psd}
\end{align}
其中$H_{\mathrm{VP}}=\nabla^2\varphi\in\R^{2\times2}$，下标VP表示变量投影；$\succeq0$表示半正定。式\eqref{eq:psd}来自系数最小二乘子问题的$H_{uu}$正定，因此$H_{zu}H_{uu}^{-1}H_{uz}$为半正定矩阵。

式\eqref{eq:schur}解释了两种更新的步长差别。位置移动引起的部分响应变化，可以由重新估计波形补偿。固定系数曲率$H_{zz}$没有扣除这部分变化，因而在耦合较强的方向上可能过大；相同梯度除以较大的正曲率，就得到较小步长。消元曲率$H_{\mathrm{VP}}$扣除了位置--波形耦合项，使更新直接针对重估波形后的残差。

在搜索区域内部的局部极小点$\bm z^*$附近，令位置误差$\bm\epsilon=\bm z-\bm z^*$。当$H_{zz}$和$H_{\mathrm{VP}}$均正定，且采用完整Newton步时，用固定系数曲率更新消元梯度的一阶误差传播为
\begin{equation}
 \bm\epsilon^{+}\approx
 \left(\bm I_2-H_{zz}^{-1}H_{\mathrm{VP}}\right)\bm\epsilon .\label{eq:localrate}
\end{equation}
式中$\bm\epsilon^+$为更新后的位置误差，Hessian取局部极小点处的值。由式\eqref{eq:psd}，迭代矩阵的特征值落在$[0,1)$；当某方向的消元曲率远小于固定系数曲率时，该方向的误差收缩因子接近1，迭代便可能出现大量小步。采用精确消元Newton曲率可去掉这一局部线性误差项。实际算法还受回溯、步长限制和Hessian近似的影响，因此用目标与导数调用次数进一步评价计算效率。

波形消元与联合精修处理两种不同的相互影响：前者允许波形随位置调整，后者允许两个位置一起调整。对联合目标$J$，双源的$4\times4$ Hessian为
\begin{equation}
 H_{\mathrm{joint}}=
 \begin{bmatrix}H_{11}&H_{12}\\H_{21}&H_{22}\end{bmatrix} .
\end{equation}
其中$H_{ij}=\partial^2J/(\partial\bm p_i\partial\bm p_j^T)\in\R^{2\times2}$。逐源精修每次只改变一个位置，联合精修同时调整两个位置，从而利用跨源项$H_{12}$和$H_{21}$所反映的相互影响。MCJ-AP采用联合L-BFGS-B更新；VP-BNOMP改进每个源内部的位置--系数耦合，不额外叠加末端联合优化。

MCJ-AP的联合精修仅接受完整目标不增的更新，即$J(\bm P^+)\le J(\bm P)$。VP-BNOMP的每次块更新降低或保持显式拟合残差，循环后的共同最小二乘重估也不增加该残差。因此，两者各自接受的目标值序列非负、单调有界，因而收敛。实际计算中，AP检查共同投影梯度；VP-BNOMP检查相邻循环的归一化残差能量变化，并设置安全迭代上限。局部搜索所得位置还取决于初始化，这也是保留多个候选的原因。

\subsection{计算量及方法关系}\label{sec:complexity}
复杂度以保留频点数$F$、站数$L$、源数$K$和初始网格点数$G=|\mathcal G|$表示。记$I$为循环精修次数，$T$为连续目标或梯度的调用总数。单次联合系数求解需形成$K\times K$ Gram矩阵$\bm A_f^H\bm A_f$，并解相应线性方程，跨频主要计算量为$\mathcal O(F(LK^2+K^3))$。固定其他源后，投影可预计算，条件网格扫描约为$\mathcal O(FGLK)$。双源全组合搜索需检查$G^2$组位置，而逐源发现每次只扫描$G$点。

\begin{table}[pos=htbp]
\centering\small
\caption{主要计算量。$I_{\mathrm B}$为稀疏贝叶斯迭代数，$S$为分段数。}\label{tab:complexity}
\begin{tabularx}{\linewidth}{@{}lX@{}}\toprule
方法或运算 & 主要量级与可观测计数\\\midrule
公共DFT & $\mathcal O(LN\log N)$；原始采样点与保留频点数\\
MUSIC／MVDR & 统计约$\mathcal O(F_sSL^2+F_sL^3)$，谱扫描$\mathcal O(F_sGL^2)$；$F_s$为分段后频点数\\
MCJ-AP & 两次条件扫描加$M$个局部位置优化；约$\mathcal O(FGLK+T_MF(LK^2+K^3))$\\
VP-BNOMP & 逐源扫描与$I$轮循环精修；约$\mathcal O(KFGL+IKT_{\mathrm{loc}}FL)$，另计共同系数重估\\
多初值／GA & 由位置目标调用总数$T$决定，约$\mathcal O(TF(LK^2+K^3))$\\
宽带SBL & 稠密字典每轮约$\mathcal O(FGL^2+FL^3)$，累计随$I_{\mathrm B}$增长\\\bottomrule
\end{tabularx}
\end{table}
表中$T_M$为全部$M$个候选分支的连续目标与梯度调用总数，$T_{\mathrm{loc}}$为一次源位置精修的调用数，$F_s$为分段DFT保留的频点数。本文通过减少实际调用次数节省计算。局部单源系数按频点独立，式\eqref{eq:schur}的$H_{uu}$为分块对角矩阵，每块仅$2\times2$；逐频计算后累加得到$2\times2$位置曲率，无需构造完整$2F\times2F$矩阵。固定系数与消元更新的单次运算均随$F$线性增长，差别主要在所需步数与循环次数。实验因此同时报告扫描点数、目标/梯度调用数和完整耗时。

AP、NOMP、BNOMP与QNOMP分别为条件搜索、残差发现、块更新及联合优化提供基础\cite{ChenHudsonYao2002,Mamandipoor2016,Yang2022,Zeng2025}。本文据此形成两条直接定位流程：MCJ-AP用多个条件起点提高恢复能力，用联合位置精修降低分支成本；VP-BNOMP保留残差发现与循环结构，通过波形消元改进局部更新。实验将候选覆盖、位置更新和波形消元分开比较，使可靠性收益与计算收益各有对应证据。
